\documentclass[12pt]{article}

%% Basic document formatting
\usepackage{amsmath, amsthm, amssymb, amsfonts}
\usepackage{mathtools}
\usepackage{mathrsfs}
\usepackage{xspace}
\usepackage{thmtools}
\usepackage{graphicx}
\usepackage{tikz-cd}
\usepackage{setspace}
\usepackage{fancyhdr}
\usepackage{titling}
\usepackage{hyperref}
\usepackage[left=0.4in,right=0.4in,top=1in,bottom=1in]{geometry}
\usepackage{float}
\usepackage{tabularx}
\usepackage{booktabs}
\usepackage[utf8]{inputenc}
\usepackage[english]{babel}
\usepackage{framed}
\usepackage[dvipsnames]{xcolor}
\usepackage{environ}
\usepackage{tcolorbox}
\tcbuselibrary{theorems,skins,breakable}

\usepackage{enumitem}
\setlist[enumerate]{leftmargin=*}
\setlist[enumerate,1]{labelindent=\parindent}
\setlist[enumerate,2]{labelindent=0pt}

% Blackboard Bold
\newcommand{\Z}{\mathbb{Z}}                 % integers
\newcommand{\Zpl}{\mathbb{Z}_{+}}           % positive integers
\newcommand{\Znn}{\mathbb{Z}_{\geq 0}}      % nonnegative integers. 
\newcommand{\Q}{\mathbb{Q}}                 % rationals
\newcommand{\Qpl}{\mathbb{Q}_{+}}           % positive Rationals
\newcommand{\R}{\mathbb{R}}                 % reals
\newcommand{\Rpl}{\mathbb{R}_{+}}           % positive Reals
\newcommand{\C}{\mathbb{C}}                 % complex numbers
\newcommand{\F}{\mathbb{F}}                 % field

% Words
\newcommand{\st}{\text{ such that }}
\newcommand{\wrt}{\text{ with respect to }}
\newcommand{\with}{\text{ with }}
\newcommand{\ie}{\text{, i.e. }}
\newcommand*{\ditto}{\text{---}\texttt{"}\text{---}}

% Operators
\newcommand{\abs}[1]{\left|#1\right|}                   % absolute value
\newcommand{\norm}[1]{\abs{\abs{#1}}}
\newcommand{\floor}[1]{\left\lfloor #1 \right\rfloor}   % floor
\newcommand{\ceil}[1]{\left\lceil #1 \right\rceil}      % ceiling
\newcommand{\powset}[1]{\mathscr{P}(#1)}

% Category Theory 
\renewcommand{\hom}{\text{Hom}}
\newcommand{\homspace}[3]{\hom_{#1}(#2, #3)}
\newcommand{\coproduct}[9]{
  \begin{tikzcd}[ampersand replacement=\&]
      \& \& #1 \arrow[dll, bend right, "#6"] \arrow[dl, "#5"] \\
   #4 \& #3 \arrow[l, "#9"] \\
      \& \& #2 \arrow[ull, bend left, "#8"] \arrow[ul, "#7"]
  \end{tikzcd}
}

\newcommand{\product}[9]{ 
  \begin{tikzcd}[ampersand replacement=\&]
      \& \& #3 \\
      #1 \arrow[r, "#7"] \arrow[urr, bend left, "#5"] \arrow[drr, bend right, "#6"] \& #2 \arrow[ur, "#8"] \arrow[dr, "#9"] \\ 
      \& \& #4
  \end{tikzcd}
}


% Algebra 
\newcommand{\Syl}[2]{\text{Syl}_{#1}{(#2)}}           % set of Sylow-p subgroups 
\newcommand{\<}{\langle}                % \<x\>, subgroup generated by x 
\renewcommand{\>}{\rangle}
\renewcommand{\index}[2]{[#1 : #2]}
\newcommand{\id}{\text{id}}             % identity element
\newcommand{\lhdneq}{%
  \mathrel{\ooalign{$\lneq$\cr\raise.22ex\hbox{$\lhd$}\cr}}}
\newcommand{\order}[1]{\text{o}(#1)}    % order of an element
\newcommand{\spec}{\operatorname{Spec}}
\newcommand{\rad}{\operatorname{rad}}   % radical of an ideal 
\newcommand{\sym}[1]{\mathfrak{S}_{#1}}
\newcommand{\aut}[1]{\text{Aut}(#1)} 
\newcommand{\gal}[2]{\text{Gal}(#1 / #2)} 
\newcommand{\inn}[1]{\text{Inn}(#1)} 
\let\oldcong\cong
\let\oldequiv\equiv
\renewcommand{\cong}{\oldequiv}
\renewcommand{\equiv}{\oldcong}

\newcommand{\triangleleftneq}{\mathrel{\ooalign{%
  $\trianglelefteq$\cr
  \hidewidth\raise0.06ex\hbox{$\not\mkern4mu$}\hidewidth\cr
}}}

\makeatletter % cycle
\newcommand{\cyc}[1]{(\mathbf{\cyc@process#1\relax})}
\def\cyc@process#1#2\relax{%
	#1%
	\ifx\relax#2\relax
	\else
		\,\cyc@process#2\relax
	\fi
}
\makeatother

\newcommand{\GL}[2]{\text{GL}_{#1}(#2)}                             % general linear group
\newcommand{\SL}[2]{\text{SL}_{#1}(#2)}                             % special linear group
\newcommand{\gl}[2]{\mathfrak{gl}_{#1}(#2)}
\renewcommand{\sl}[2]{\mathfrak{sl}_{#1}(#2)}
\newcommand{\so}[2]{\mathfrak{so}_{#1}(#2)}
\newcommand{\su}[1]{\mathfrak{su}(#1)}

% Linear Algebra
\newcommand{\bmat}[1]{\begin{bmatrix}#1\end{bmatrix}}   % bracketed matrix
\newcommand{\rank}{\operatorname{rank}}                 % rank
\newcommand{\nullity}{\operatorname{nullity}}           % nullity
\newcommand{\tr}{\operatorname{tr}}

% Combinatorics
\newcommand{\qnum}[1]{\left[#1\right]_q}                % q-analog
\newcommand{\qfact}[1]{\left[#1\right]!_q}              % q-analog factorial 
\newcommand{\qbinom}[2]{\binom{#1}{#2}_q}               % Gaussian binomial coefficient

% Topology/Analysis
\newcommand{\ball}[2]{\text{B}_{#1}(#2)}  % B_r(x): open r-balls around x
\newcommand{\diam}{\text{diam}}           % diamter of a set in metric space 
\newcommand{\lowint}[2]{\underline{\int_{#1}^{#2}}}
\newcommand{\upint}[2]{\overline{\int}_{#1}^{#2}}

% Misc Notation
\newcommand{\oneton}[1]{\{1, \ldots, #1\}}
\newcommand{\defeq}{\vcentcolon=}   % :=
\newcommand{\eqdef}{=\vcentcolon}   % =: 
\renewcommand{\bf}[1]{\textbf{#1}}
\newcommand{\im}{\operatorname{im}}
\newcommand{\fnrest}[2]{\left.#1\right|_{#2}}
\newcommand{\isom}{\overset{\sim}{\rightarrow}} 


\renewcommand{\thesection}{\arabic{section}.}
\renewcommand{\thesubsection}{(\alph{subsection})}

\newtheorem{claim}{Claim}
\newtheorem*{lemma}{Lemma}

%% Headers & title setup
\newcommand{\course}{Math100C}
\newcommand{\myname}{Jay Ser}
\setlength{\headheight}{14.5pt}
\pagestyle{fancy}
\fancyhf{}
\renewcommand{\headrulewidth}{0.4pt}
\lhead{\course}
\rhead{\myname}
\cfoot{\thepage}
\setlength{\droptitle}{-4em} 
\title{\course\ - WK3} 
\author{\myname}
\date{2026.04.17}

\begin{document}
\maketitle
\thispagestyle{fancy}

%------------------------------------------------------------------------------%

\section{An Element of $\mathrm{GL}_n(\mathbb{C})$ of Infinite Order}

Let $n$ be a positive integer and let $A$ be the $n \times n$ matrix over
$\mathbb{C}$ such that $A_{ij} = 1$ if $j = i$ or $j = i+1$ and $0$
otherwise.  For example, for $n = 4$ we have
\[
  A =
  \begin{pmatrix}
    1 & 1 & 0 & 0 \\
    0 & 1 & 1 & 0 \\
    0 & 0 & 1 & 1 \\
    0 & 0 & 0 & 1
  \end{pmatrix}.
\]
Prove that $A$ is an element of $\mathrm{GL}_n(\mathbb{C})$ of infinite order.

\emph{Hint:} look at the $(1,2)$-entries of the powers of $A$.

\emph{Optional second part (only do this if you are comfortable with the
underlying linear algebra):} use this to give another proof that every matrix
in $\mathrm{GL}_n(\mathbb{C})$ of finite order is diagonalizable, using the
fact that every matrix over $\mathbb{C}$ can be conjugated into Jordan normal
form.

\noindent\dotfill

Pick an arbitrary $n$ and let $A$ be the $n \times n$ matrix as described. 
To show that $A$ has infinite order in $\GL_n{\C}$, it suffices to show that the $1, 2$th entry $A^k$ is $k$. 
Induct on $k$; 
For the sake of calculation, I further show that for each $k$, the $2, 2$th entry is 1. 
This is true by definition of $A$ when $k = 1$. 
Suppose the $1, 2$th entry of $A^{k - 1}$ is $k - 1$ and the $2, 2$th entry of $A^{k - 1}$ is 1. 
Let $A = [a_{ij}]$ and $A^{k - 1} = [a_{ij}^{(k - 1)}]$. 
Recalling that upper triangle matrices form a subgroup of $\GL_{n}(\C)$, one obtains the following simple calculation for $a_{12}^{(k)}$: 
\begin{align*}
  a_{12}^{(k)} & = \sum_{t = 1}^{n} a_{1t} a_{t2}^{(k - 1)} \\
               & = a_{11} a_{12}^{(k - 1)} + a_{12} a_{22}^{(k - 1)} \\ 
               & = k - 1 + 1 = k \\ 
    a_{22}^{(k)} & = \sum_{t = 1}^{n} a_{2t} a_{t2}^{(k - 1)} \\ 
                 & = a_{22}a_{22}^{(k - 1)} + a_{23} a_{32}^{(k - 1)} \\ 
                 & = 1 + 0 = 1,
\end{align*}
as was to be shown. 

%------------------------------------------------------------------------------%

\section{A Character with $\langle \chi, \chi \rangle = 2$}

Let $\chi$ be a character of a finite group $G$ such that
$\langle \chi, \chi \rangle = 2$.  Using orthogonality, prove that the
corresponding representation is the direct sum of two distinct irreducible
representations.

\noindent\dotfill

Since $\<\chi, \chi\> \neq 1$, $\chi$ is a reducible character. 
Let 
$$\chi = \sum_{i = 1}^{n} c_i \chi_i $$
where each $\chi_{i}$ is irreducible. 
Then 
\begin{align*}
  \<\chi, \chi\> & = \<\sum_{i = 1}^{n} c_{i} \chi_{i}, \sum_{i = 1}^{n} c_{i} \chi_{i}\> \\ 
                 & = \sum_{1 \leq i, j \leq n} c_i c_j \<\chi_{i}, \chi_{j}\> \\ 
                 & = \sum_{i = 1}^{n} c_{i}^{2} 
\end{align*}
where the third equality follows from the orthogonality of irredicuble characters. 
Because $\<\chi, \chi\> = 2$, it follows from simple arithmetic that $n = 2$ and $c_{1} = c_{2} = 1$, i.e., 
$$\chi = \chi_{1} + \chi_{2}.$$

%------------------------------------------------------------------------------%

\section{Products of Characters}

Let $\chi_1, \chi_2$ be two characters of a finite group $G$.

\subsection{$\chi_1 \chi_2$ is a Character}

Prove that $\chi_1 \chi_2$ (that is, the function taking $g$ to
$\chi_1(g)\chi_2(g)$) is also a character of $G$.

\emph{Hint:} during the proof of orthogonality of characters, we will see that
$\overline{\chi_1}\chi_2$ is a character of $G$.  Use that same logic but
starting from $\overline{\chi_1}$ in place of $\chi_1$.

\noindent\dotfill

If $R_1: G \rightarrow \GL(V)$ is the representation corresponding to $\chi_1$, then 
$$R_{1}':G \rightarrow \GL(V); \hspace{1.5em} g \mapsto \overline{R_g}$$ 
i.e., postcomposing $R$ with complex conjugation of a linear operator on $V_1$, is another group representation of $G$.
Evidently, $\chi_{1}' = \overline{\chi_1}$. 
Hence $\chi_1 \chi_2$ is a character of $G$. 


\subsection{Irreducibility When $\chi_1(1) = 1$}

Check that if $\chi_1(1) = 1$ and $\chi_2$ is irreducible, then
$\chi_1 \chi_2$ is irreducible.

\noindent\dotfill

Since $\chi_1$ is a one-dimensional representation, $\overline{\chi_{1}(g) \chi_{1}(g)} = 1$ for all $g \in G$ because $\chi_{1}(g)$ is a root of unity. 
It follows 
\begin{align*}
  \<\chi_1 \chi_2, \chi_1 \chi_2\> 
  & = \frac{1}{\#G} \sum_{g \in G} \overline{\chi_1 \chi_2} \chi_1 \chi_2 \\ 
  & = \frac{1}{\#G} \sum_{g \in G} \overline{\chi_1} \chi_1 \overline{\chi_2} \chi_2 \\ 
  & = \frac{1}{\#G} \sum_{g \in G} \overline{\chi_2} \chi_2 = \<\chi_2, \chi_2\> = 1,
\end{align*}
hence $\chi_1 \chi_2$ is irreducible. 

\subsection{Failure When $\chi_1$ is Merely Irreducible}

Check that the previous point fails if we replace the condition
$\chi_1(1) = 1$ by the condition that $\chi_1$ is irreducible.

\emph{Hint:} look at $\sym_3$.

\noindent\dotfill

Recall 
\[
  \begin{array}{c|ccccc}
    & 1   & \cyc{12}  & \cyc{123}   \\
    \hline
    \chi & 2 & 0 & -1 
  \end{array}
\]
is the only irreducible character of $\sym_3$ with dimension greater than 1. 
Clearly, $\chi^2$ is not an irreducible character of $\sym_3$. 

%------------------------------------------------------------------------------%

\section{Character Table of $\sym_4$}

Build the character table for $\sym_4$ using the following steps.

\noindent \dotfill 

First, let's determine the conjugacy classes of $\sym_{4}$. 
The conjugacy classes in $\sym_4$ are just the possible cycle types: 
$$(1), \, (12), \, (12)(34), \, (123), \, (1234)$$
give the representatives of all the conjugacy classes in $\sym_4$. 
By doing a bit of combinatorics, one can check there are 
\begin{itemize}
  \item One identity 
  \item $\binom{4}{2} = 6$ transpositions 
  \item $\binom{4}{2} / 2 = 3$ double transpositions 
  \item $\binom{4}{3} \cdot 2 = 8$ three-cycles 
  \item $3! = 6$ four-cycles
\end{itemize}
which shows that $\sym_4$ has five irreducible characters. 

\subsection{Trivial and Sign Representations}
As with $S_3$, start with the trivial irreducible representation and the
nontrivial one-dimensional representation coming from signs of permutations.

\noindent\dotfill

\[
  \begin{array}{c|ccccc}
    & (1) & (6) & (3) & (8) & (6) \\
    & 1   & \cyc{12}   & \cyc{12}\cyc{23}   & \cyc{123}   & \cyc{1234}   \\
    \hline
    \chi_1 & 1 & 1 &  1 &  1 &  1 \\
    \chi_2 & 1 & -1 & 1 & 1 &  -1 \\
    \chi_3 &  \\
    \chi_4 &  \\ 
    \chi_5 & \\ 
  \end{array}
\]

\subsection{Standard Three-Dimensional Representation}

Again as with $\sym_3$, take the standard four-dimensional representation, then
subtract a copy of the trivial representation, and use orthogonality to check
that the result is irreducible.

\noindent\dotfill

Let $\chi_{R}$ be the character of the representation that sends a permutation to its permutation matrix. 
Given any $w \in \sym_n$, recall that the trace of the permutation matrix is precisely the number of points fixed by $w$. 
\[
  \begin{array}{c|ccccc}
    & 1   & \cyc{12}   & \cyc{12}\cyc{23}   & \cyc{123}   & \cyc{1234}   \\
    \hline
    \chi_R & 4 & 2 & 0 &  1 &  0 \\
  \end{array}
\]
One can use this table to calculate 
$$\<\chi_R, \chi_R\> = \frac{1}{24}\sum_{w \in \sym_4}\overline{\chi_{R}(w)}\chi_{R}(w) = 48/24 = 2$$
which means $\chi_{R}$ is the sum of two irreducible characters. 
As usual, the vector $[1, 1, 1, 1]^{t}$ is fixed by the permutation matrices, so one said irreducible character is the trivial character.  
Hence $\chi_{3} \defeq \chi_R - \chi_1$ is the standard three-dimensional character. 
\[
  \begin{array}{c|ccccc}
    & (1) & (6) & (3) & (8) & (6) \\
    & 1   & \cyc{12}   & \cyc{12}\cyc{23}   & \cyc{123}   & \cyc{1234}   \\
    \hline
    \chi_1 & 1 & 1 &  1 &  1 &  1 \\
    \chi_2 & 1 & -1 & 1 & 1 &  -1 \\
    \chi_3 & 3 & 1 & -1 & 0 & -1 \\
    \chi_4 &  \\ 
    \chi_5 & \\ 
  \end{array}
\]

\subsection{A Second Three-Dimensional Representation}

Add a second three-dimensional representation by multiplying by the nontrivial
one-dimensional representation.

\noindent\dotfill

Let $\chi_{4} \defeq \overline{\chi_2}\chi_3$ be another character of $G$. 
\[
  \begin{array}{c|ccccc}
    & 1   & \cyc{12}   & \cyc{12}\cyc{23}   & \cyc{123}   & \cyc{1234}   \\
    \hline
    \chi_4 & 3 & -1 & -1 &  0 &  1 \\
  \end{array}
\]
One can calculate $\<\chi_4, \chi_4\> = 1$, hence $\chi_4$ is an irreducible character of $G$. 
\[
  \begin{array}{c|ccccc}
    & (1) & (6) & (3) & (8) & (6) \\
    & 1   & \cyc{12}   & \cyc{12}\cyc{23}   & \cyc{123}   & \cyc{1234}   \\
    \hline
    \chi_1 & 1 & 1 &  1 &  1 &  1 \\
    \chi_2 & 1 & -1 & 1 & 1 &  -1 \\
    \chi_3 & 3 & 1 & -1 & 0 & -1 \\
    \chi_4 & 3 & -1 & -1 & 0 & 1 \\ 
    \chi_5 & \\ 
  \end{array}
\]

\subsection{The Remaining Two-Dimensional Character}

Note that the last character must be two-dimensional.  Use orthogonality to
fill in its values.

\noindent\dotfill

By the identity $24 = \#\sym_4 = \chi_{1}(1)^2 + \ldots + \chi_{5}(1)^{2}$, obtain $\chi_5(1) = 2$. 
The columns of the character table must be orthogonal, forcing the rest of the entries: 
\[
  \begin{array}{c|ccccc}
    & (1) & (6) & (3) & (8) & (6) \\
    & 1   & \cyc{12}   & \cyc{12}\cyc{23}   & \cyc{123}   & \cyc{1234}   \\
    \hline
    \chi_1 & 1 & 1 &  1 &  1 &  1 \\
    \chi_2 & 1 & -1 & 1 & 1 &  -1 \\
    \chi_3 & 3 & 1 & -1 & 0 & -1 \\
    \chi_4 & 3 & -1 & -1 & 0 & 1 \\ 
    \chi_5 & 2 & 0 & 2 & -1 & 0\\ 
  \end{array}
\]

%------------------------------------------------------------------------------%

\section{Partial Character Table of a Group of Order 20}

The following is part of the character table of a group of order $20$.  Fill
in the last row, then determine the orders of the elements $a, b, c, d$.
(Here $i$ has its usual meaning of a square root of $-1$.)
\[
  \begin{array}{c|ccccc}
    & (1) & (4) & (5) & (5) & (5) \\
    & 1   & a   & b   & c   & d   \\
    \hline
    \chi_1 & 1 & 1 &  1 &  1 &  1 \\
    \chi_2 & 1 & 1 & -1 & -1 &  1 \\
    \chi_3 & 1 & 1 & -i &  i & -1 \\
    \chi_4 & 1 & 1 &  i & -i & -1 \\
    \chi_5 &   &   &    &    &
  \end{array}
\]

\noindent\dotfill

Since $\sum_{i = 1}^{5} \chi_{i}(1)^{2} = 20$, 
$$\chi_{5}(1) = 4.$$
The columns of a character table are orthogonal, so the following table is forced after knowing the value of $\chi_{5}(1)$: 
\[
  \begin{array}{c|ccccc}
    & (1) & (4) & (5) & (5) & (5) \\
    & 1   & a   & b   & c   & d   \\
    \hline
    \chi_1 & 1 & 1 &  1 &  1 &  1 \\
    \chi_2 & 1 & 1 & -1 & -1 &  1 \\
    \chi_3 & 1 & 1 & -i &  i & -1 \\
    \chi_4 & 1 & 1 &  i & -i & -1 \\
    \chi_5 & 4 & -1 &  0 & 0 & 0
  \end{array}
\]

Next, let's determine the orders of $a, b, c$, and $d$. 
Recall that conjugates have the same order in the group. 
By Cauchy's theorem, there is an element of order 2 and an element of order 5. 
Let $x \in G$ be an order 5 element. 
Then the value of the one-dimensional characters at $x$ must all be fifth roots of unity. 
Only the column for $a$ satisfies this condition, so $x$ must be conjugate to $a$. 
Similarly, if $y \in G$ is an order 2 element, the one-dimensional characters should evaluate at $y$ to either $1$ or $-1$.
Of the remaining columns, only the column for $d$ satisfies this condition. 
\begin{align*}
  o(a) & = 5 \\ 
  o(d) & = 2
\end{align*}
What about $b$ and $c$? 
Well, since the one-dimensional characters evaluate at $b$ and $c$ to a fourth root of unity, the order of $b$ and $c$ must be divisible by $4$. 
Furthermore, the order of $b$ and $c$ divides the group order, which is 20. 
So $o(b), o(c) \in \{4, 20\}$. 
If either element has order 20, then $G$ is cyclic. 
This contradicts the assumption given by the character table that the group is nonabelian; 
if the group was abelian, then each of the 20 elements of $G$ would form its own conjugacy class. 
Hence, 
\begin{align*}
  o(b) & = 4 \\ 
  o(c) & = 4
\end{align*}

%------------------------------------------------------------------------------%

\section{Non-Simplicity via Characters, and Simplicity of $A_5$}

Let $G$ be a finite group.  Show that if $G$ is not simple, then we can find a
nontrivial irreducible character $\chi$ of $G$ and an element $g \neq 1$ in
$G$ such that $\chi(g) = \chi(1)$.  Then apply this criterion to the character
table of $A_5$ to see that $A_5$ is a simple group.

\noindent\dotfill

Suppose $G$ is not simple; 
take a proper nontrivial normal subgroup $N$ of $G$. 
Since $G / N$ is not the trivial group, there is a nontrivial irreducible representation $\tilde{R}: G / N \rightarrow \GL(V)$. 
Let 
$$R: G \rightarrow \GL(V); \hspace{1.5em} R = \tilde{R} \pi.$$ 
Since a composition of homomorphisms is a homomorphism, $R$ is a group representation. 
Namely, it is an irreducible representation because $\tilde{R}$ is irreducible. 
Hence for any $n \neq 1 \in N$, $R(n) = \tilde{R}(N) = R(1)$,
i.e., 
we have found an irreducible character of a group $G$ such that 
$$\exists n \neq 1 \in G \st \chi(n) = \chi(1).$$

Recall that the character table for $A_5$ looks like 
\[
  \begin{array}{c|ccccc}
    & (1) & (15) & (20) & (12) & (12) \\
    & \id   & \cyc{12}{34}   & \cyc{123}   & \cyc{12345}   & \cyc{12354}   \\
    \hline
    \chi_1 & 1 & 1 &  1 &  1 &  1 \\
    \chi_2 & 3 & 0 & -1 & \bar{\alpha} &  \alpha \\
    \chi_3 & 3 & 0 & -1 &  \alpha & \bar{\alpha} \\
    \chi_4 & 4 & 1 &  0 & -1 & -1 \\
    \chi_5 & 5 & -1 &  1 & 0 & 0
  \end{array}
\]
where $\alpha$ is a complex number with nonzero complex part and $\bar{\alpha}$ denotes its complex conjugate. 
There is no nontrivial irreducible character $\chi$ of $A_5$ such that $\exists w \neq 1 \in A_5$ with $\chi(w) = \chi(1)$. 
One concludes that $A_5$ is simple. 

%------------------------------------------------------------------------------%

\section{Permutation Representation of $\sym_3$ Acting on Itself by Conjugation}

Let $R$ be the permutation representation coming from the action of $\sym_3$ on
itself via conjugation.  Decompose the associated character into irreducibles,
then verify that $R$ is \emph{not} isomorphic to the regular representation of
$\sym_3$.

\noindent\dotfill

Let the following be the associated permutation representation of this action: 
\begin{align*}
  \varphi: \sym_{3} \rightarrow \sym_{6}; \hspace{1.5em} & w \mapsto \sigma_{w} \\ 
  \sigma_{w}: \sym_{3} \rightarrow \sym_{3}; \hspace{1.5em} & z \mapsto wzw^{-1}
\end{align*}

The following gives the action of $\cyc{12}$ on $\sym_{3}$: 
\begin{align*}
  \sigma_{\cyc{12}}\id = \id; \hspace{1.5em} & \sigma_{\cyc{12}}\cyc{12} = \cyc{12} \\ 
  \sigma_{\cyc{12}}\cyc{123} = \cyc{132}; \hspace{1.5em} & \sigma_{\cyc{12}}\cyc{13} = \cyc{23} \\ 
  \sigma_{\cyc{12}}\cyc{132} = \cyc{123}; \hspace{1.5em} & \sigma_{\cyc{12}}\cyc{23} = \cyc{13}
\end{align*}
If the elements of $\sym_{3}$ are ordered as $\{ \id, \cyc{12}, \cyc{13}, \cyc{23}, \cyc{123}, \cyc{132}\}$, then $\sigma_{\cyc{12}}$ can be identified as $\cyc{34}\cyc{56}$ in $\sym_{6}$. 
Via a similar analysis, $\sigma_{\cyc{123}}$ can be identified as $\cyc{243}$ in $\sym_{6}$. 
Realizing this permutation representation as a matrix representation $R$ via permutation matrices gives the following character $\chi$ of $R$: 
\[
  \begin{array}{c|ccc}
    & 1   & \cyc{12}   & \cyc{123} \\
    \hline
    \chi & 6 & 2 & 3
  \end{array}
\]

Recall that below is the character table for $\sym_{3}$: 
\[
  \begin{array}{c|ccc}
    & 1   & \cyc{12}   & \cyc{123} \\
    \hline
    \chi_1 & 1 & 1 & 1 \\ 
    \chi_2 & 1 & -1 & 1 \\ 
    \chi_3 & 2 & 0 & -1
  \end{array}
\]
$\chi$ can be decomposed into irreducibles by projecting $\chi$ onto each of the irreducible characters: 
\begin{align*}
  \<\chi, \chi_{1}\> & = \frac{1}{6}(6 + 6 + 6) = 3 \\ 
  \<\chi, \chi_{2}\> & = \frac{1}{6}(6 -6 + 6) = 1 \\ 
  \<\chi, \chi_{3}\> & = \frac{1}{6}(12 + 0 - 6) = 1
\end{align*}
Hence $\chi = 3\chi_1 + \chi_2 + \chi_1$. 
The decomposition of the regular representation is $\chi_{R} = \chi_1 + \chi_2 + 2\chi_3$, so one concludes that $\chi$ and $\chi_{R}$ are not isomorphic. 

%------------------------------------------------------------------------------%

\section{Converse to Schur's Lemma}

Prove the following converse to Schur's lemma: if $R \colon G \to \mathrm{GL}(V)$
is a vector space representation of a finite group $G$ and the only
$G$-invariant linear operators on $V$ are multiplication by scalars, then $R$
is irreducible.

\emph{Hint:} use Maschke's theorem.

\noindent\dotfill

Suppose $R: G \rightarrow \GL(V)$ is reducible. 
Then there is a nontrivial $G$-invariant proper subspace $W \subsetneq V$; 
by Maschke's theorem, $V = W \oplus W^{\perp}$ with $W^{\perp}$ $G$-invariant as well. 
Let 
$$\pi: V \rightarrow V; \: v \mapsto w\hspace{1.5em} (v = w + w' \mid w \in W, w' \in W^{\perp})$$
be the projection onto $W$. 
For the rest of the paragraph, let $w$ always denote a vector in $W$ and $w'$ a vector in $W^{\perp}$. 
Projections are clearly linear operators. 
Furthermore, $\pi$ is a $G$-invariant linear operator: 
for any $g \in G$ and $v = w + w' \in V$, 
\begin{align*}
  R_{g} \pi R_{g^{-1}}(v) & = R_{g} \pi(R_{g^{-1}}(w) R_{g^{-1}}(w')) \\ 
                          & = R_{g} R_{g^{-1}}(w) \\ 
                          & = w = \pi(v)
\end{align*}
where the second equality follows because of the $G$-invariance of $W$ and $W^{\perp}$; 
$R_{g^{-1}}(w) \in W$ and $R_{g^{-1}}(w') \in W^{\perp}$. 
Finally, one can check that $\pi$ is not a scalar linear operator. 
If it was, let $\pi(v) = \mu v$ for all $v = w + w' \in V$. 
Then 
$$w = \pi(v) = \mu(w + w') \iff (\mu - 1)w + \mu w' = 0$$ 
and because $w$ and $w'$ are linearly independent, $\mu = 0$. 
$\mu = 0$ says $\pi$ is the 0 map, which contradicts the definition of $\pi$ as the projection onto $W$ and the assumption that $W$ is an irreducible subspace of $V$, that is, $W \neq 0$. 
This proves the contrapositive to the desired statement. 

\end{document}
